\documentclass[a4paper,fleqn]{cas-sc}

% References and mathematics
\usepackage[numbers,sort&compress]{natbib}
\usepackage{amsmath,amssymb}
\usepackage{xurl} % Allow long repository URLs to wrap within the text block.

% Figures
\usepackage{graphicx}
\usepackage[export]{adjustbox}
\usepackage{float}
\usepackage{placeins}
\usepackage{capt-of}

% Tables
\usepackage{booktabs}
\usepackage{array}
\usepackage{tabularx}
\usepackage{multirow}

% Algorithms
\usepackage{algorithm}
\usepackage{algpseudocode}

% Figure paths
\graphicspath{{figures/}{./figures/}{./}}
\DeclareGraphicsExtensions{.pdf,.png,.jpg,.jpeg}

% Table column types
\newcolumntype{Y}{>{\raggedright\arraybackslash}X}
\newcolumntype{Z}[1]{>{\centering\arraybackslash}p{#1}}

% A compact non-floating figure command used throughout the paper.
% Usage:
% \paperfig[trim={...},clip]{0.92\linewidth}{fig.pdf}{Caption.}{fig:label}
\newcommand{\paperfig}[5][]{%
  \begin{center}
  \begin{minipage}{\linewidth}
  \centering
  \includegraphics[
    width=#2,
    max height=0.52\textheight,
    keepaspectratio,
    #1
  ]{#3}%
  \vspace{-0.35em}
  \captionof{figure}{#4}
  \label{#5}
  \end{minipage}
  \vspace{-0.65em}
  \end{center}
}

\let\printorcid\relax

\begin{document}

\setlength{\abovedisplayskip}{0pt plus 1pt minus 0pt}
\setlength{\belowdisplayskip}{5pt plus 1pt minus 1pt}
\setlength{\abovedisplayshortskip}{0pt plus 1pt minus 0pt}
\setlength{\belowdisplayshortskip}{5pt plus 1pt minus 1pt}

\shorttitle{Few-Shot DT Instantiation}

\title[mode=title]{Adaptation-Aware Model Selection for Few-Shot Digital Twin Instantiation}

\begin{abstract}
As an important modeling tool for managing computing centers, digital twin
(DT) platforms reuse pre-trained DT models from existing centers to reduce the
cost of DT instantiation in large-scale computing networks. Target-center
heterogeneity in resource configurations and workload characteristics
complicates direct model reuse when target adaptation uses a small support set.
However, existing few-shot adaptation methods typically retain a predefined
source architecture and therefore do not directly address target-specific
model selection for heterogeneous centers under per-model deployment limits.
To address this challenge, we propose Model Selection and Adaptation for
Digital Twin Instantiation (MSA-DTI), an adaptation-aware model-selection
algorithm for DT instantiation with few-shot target adaptation under deployment
limits. MSA-DTI constructs a reusable candidate bank of source-trained models
with diverse architectures, filters candidates that violate the operation-count
and parameter-count limits, and adapts the remaining candidates using the same
target support set and update budget. It then compares post-adaptation
performance on a separate labeled validation set and replaces a predefined reference model only when the
best-performing alternative achieves a calibrated relative reduction in
validation loss. On the controlled synthetic benchmark, MSA-DTI achieves the lowest mean
squared error (MSE) among the evaluated methods and reduces case-wise MSE by
14.60\% on average relative to the pre-training-and-fine-tuning baseline,
while satisfying both deployment limits. In a separately source-trained and calibrated Bitbrains real-workload
evaluation, we observe a 32.70\% mean case-relative test-MSE reduction
versus PT+FT (95\% VM-cluster CI [7.67\%,51.25\%]), with increased
target-side adaptation cost.
\end{abstract}

\begin{keywords}
digital twin
\sep few-shot adaptation
\sep model instantiation
\sep architecture selection
\sep model deployment limits
\sep heterogeneous computing centers
\end{keywords}

\maketitle


\section{Introduction}
\label{sec:introduction}

As an important modeling tool for managing computing centers, a shared digital
twin (DT) platform maintains center-specific models for monitoring, prediction,
analysis, and decision making \cite{ref1,ref2,ref3,ref4}. In large-scale computing
networks, reusing and adapting pre-trained DT models instead of constructing
each target model from scratch reduces DT-instantiation cost
\cite{ref8,ref11}. Each pre-trained model comprises an architecture and
source-trained parameters. The architecture determines the model structure and
deployment-related complexity, whereas the parameters provide the
initialization for target adaptation. For a new or reconfigured center, differences in resource configurations,
monitored-variable availability, workload patterns, and data distributions
complicate direct reuse of a source model, while target adaptation uses a small
support set. The source architecture is therefore not necessarily suitable for
the target center. In addition, large-scale DT platforms impose per-model
resource budgets, represented here by limits on estimated operation count per
inference and parameter count. Target-center instantiation therefore requires
selecting a suitable reusable model and adapting its parameters from the small
support set while satisfying both deployment limits. Figure~\ref{fig:scenario}
illustrates this setting.

Existing few-shot adaptation methods reuse pre-trained models
\cite{ref9,ref10} and have also been applied to DT modeling \cite{ref11}.
However, these methods typically retain the source architecture and therefore
do not directly address target-specific model selection for heterogeneous target
centers under deployment limits. Deployment-aware architecture-selection
methods consider model performance together with deployment cost
\cite{ref12,ref13,ref26}, but do not directly address few-shot adaptation of
reusable pre-trained models. Methods that jointly consider architecture
decisions and few-shot adaptation \cite{ref30,ref31,ref32} typically rely on
target-side architecture or adaptation-structure search. The remaining
challenge is therefore to select a suitable reusable pre-trained DT model under
per-model deployment limits and adapt its parameters from a small target
support set.

\paperfig[trim={5bp 15bp 5bp 5bp},clip]{0.8\linewidth}
{fig1.pdf}
{System setting for target-specific DT instantiation.}
{fig:scenario}

To address this challenge, we propose Model Selection and Adaptation for
Digital Twin Instantiation (MSA-DTI), an adaptation-aware model-selection
algorithm for DT instantiation with few-shot target adaptation under deployment
limits. MSA-DTI first
constructs a reusable candidate bank of source-trained models with diverse
architectures. For a target center, it filters candidates that violate the
deployment limits and adapts the remaining feasible candidates under a common adaptation
budget. It then compares the adapted candidates using validation performance
and retains a predefined reference model unless the best-performing alternative
achieves a sufficient reduction in validation loss. Thus, model selection is
based on post-adaptation performance without target-side architecture or
adaptation-structure search.

The main contributions are summarized below.

\begin{itemize}

\item We formulate DT instantiation as jointly selecting among reusable
source-trained models with different architectures and adapting the selected
model's parameters from a small target support set, subject to hard per-model
operation-count and parameter-count limits.

\item We propose a feasibility-first, adaptation-aware model-selection strategy
that filters out infeasible candidates before target adaptation and compares
the remaining candidates only after adapting them using the same target support
set and update budget.

\item We introduce a reference-model replacement rule with a relative-improvement
threshold calibrated on disjoint development data. An adapted alternative
replaces the predefined reference only when it achieves the calibrated relative
reduction in validation loss.

\end{itemize}

The remainder of this paper is organized as follows.
Section~\ref{sec:related} reviews related work.
Section~\ref{sec:problem} defines the system model and problem.
Section~\ref{sec:method} presents MSA-DTI.
Section~\ref{sec:experiments} reports the experiments, and
Section~\ref{sec:conclusion} concludes the paper.

\section{Related Work}
\label{sec:related}

\subsection{Digital Twin Modeling and Platform Deployment}

In communication and computing systems, digital twins represent system states
and services while supporting model updates \cite{ref5,ref6,ref7,ref8}.
At the deployment level, related work considers service orchestration and
computation offloading under resource constraints \cite{ref27,ref28}.
However, DT modeling studies typically assume a predefined model structure,
whereas platform-level research focuses on resource decisions. Neither line of work directly addresses architecture selection together with
few-shot target adaptation under deployment limits.

\subsection{Limited-Data Model Adaptation}

Limited-data methods leverage source or prior knowledge for target-model
adaptation \cite{ref9,ref10,ref24}. Recent approaches improve cross-task and
cross-domain adaptation through context modeling and transferable
representations \cite{ref20,ref21,ref22,ref23}. Limited-data learning has also
been applied to DT construction and workload prediction \cite{ref11,ref25}.
Most parameter-adaptation methods operate with a fixed architecture. For
heterogeneous target centers, a fixed architecture is not necessarily suitable,
motivating selection among reusable source-trained models.

A direct neighboring line of work selects pre-trained models from a finite
model hub for a target task. LogME assesses model suitability before
target fine-tuning and supports regression, whereas model-hub ranking and
tuning further study how pre-trained candidates can be exploited
\cite{ref35,ref36}. In contrast, the present study combines two hard
per-model deployment limits with common-budget few-shot adaptation of
feasible source-trained candidates, post-adaptation validation comparison,
and a calibrated option to retain an explicit reference. The constituent
ideas of pre-trained model ranking and model reuse are not claimed as new.

Model selection introduces a related challenge because candidate rankings
based on validation estimates are subject to uncertainty
\cite{ref33,ref34}. This motivates a conservative decision rule when adapted
reusable models are compared on target validation data.

\subsection{Deployment-Aware Architecture Selection}

Deployment-aware neural architecture search evaluates model performance
together with deployment cost \cite{ref12,ref13,ref26}. Few-shot, zero-shot,
and training-free NAS variants reduce the cost of architecture evaluation
during search \cite{ref14,ref15,ref16}. In Ref.~\cite{ref14}, ``few-shot''
denotes evaluation with several sub-supernets rather than adaptation from a
small target support set.

Among meta-learning-based NAS methods, T-NAS learns transferable architecture
knowledge for new-task adaptation \cite{ref31}. H-Meta-NAS further incorporates
hardware constraints into task-adaptive search \cite{ref30}. Neural
fine-tuning search instead starts from a pre-trained backbone and searches over
adaptation structures, including which components to fine-tune or augment
\cite{ref32}. Together, these methods extend task adaptation beyond parameter
updates to architecture or adaptation-structure decisions.

To address target-specific model selection under deployment limits without
target-stage search for new architectures or adaptation structures, MSA-DTI
selects among a finite bank of separately source-trained reusable models using
deployment feasibility and post-adaptation performance.

\section{System Model and Problem Formulation}
\label{sec:problem}

\subsection{System Setting and Target Model}

We consider a shared DT platform that maintains center-specific models for
heterogeneous computing centers.

Let $c$ denote a new or reconfigured target center. To enable model reuse
across centers with different monitored-variable availability, all observations
follow a common metric schema, with unavailable variables represented by masks
rather than omitted. Let $\mathcal{X}$ denote the common input space defined by
this metric schema, and let $\mathcal{Y}$ denote the output space of the DT
task. Let $\mathcal{A}$ denote the set of model architectures available to the
shared DT platform. For $a\in\mathcal{A}$, let $\Theta(a)$ denote the
trainable-parameter space associated with architecture $a$, and let
$\theta\in\Theta(a)$ denote a particular parameter setting. A model is
represented as

\begin{equation}
f_{a,\theta}:\mathcal{X}\rightarrow\mathcal{Y},
\qquad
a\in\mathcal{A},\;
\theta\in\Theta(a).
\label{eq:model-family}
\end{equation}

The DT model instantiated for center $c$ is represented as

\begin{equation}
\mathrm{DT}_{c}
=
\left(
a_c,\theta_c
\right),
\qquad
a_c\in\mathcal{A},\;
\theta_c\in\Theta(a_c),
\label{eq:dt-instance}
\end{equation}
where \(a_c\) and \(\theta_c\) denote the target architecture and target-specific
parameters, respectively. During instantiation, the available target data are
represented by a support set and a validation set.

\begin{equation}
\begin{aligned}
\mathcal{D}_{c}^{\mathrm{sup}}
&=
\left\{
(x_{c,i}^{\mathrm{sup}},y_{c,i}^{\mathrm{sup}})
\right\}_{i=1}^{K_c},
\\
\mathcal{D}_{c}^{\mathrm{val}}
&=
\left\{
(x_{c,j}^{\mathrm{val}},y_{c,j}^{\mathrm{val}})
\right\}_{j=1}^{V_c},
\\
\mathcal{D}_{c}^{\mathrm{tar}}
&=
\mathcal{D}_{c}^{\mathrm{sup}}
\cup
\mathcal{D}_{c}^{\mathrm{val}},
\qquad
\mathcal{D}_{c}^{\mathrm{sup}}
\cap
\mathcal{D}_{c}^{\mathrm{val}}
=
\varnothing .
\end{aligned}
\label{eq:target-data}
\end{equation}

In these sets,
$(x_{c,i}^{\mathrm{sup}},y_{c,i}^{\mathrm{sup}})$ and
$(x_{c,j}^{\mathrm{val}},y_{c,j}^{\mathrm{val}})$ denote the $i$th support
and $j$th validation input-output pairs, respectively. Their union
$\mathcal{D}_{c}^{\mathrm{tar}}$ denotes the target data available during
instantiation. $K_c$ denotes the number of target support samples used for
parameter adaptation, whereas $V_c$ denotes the number of validation samples
used for post-adaptation candidate comparison. The present study concerns
determining $(a_c,\theta_c)$ during target-center instantiation. Subsequent
model updates are outside its scope.

\subsection{Reusable Models and Deployment Limits}

The shared DT platform maintains source-trained models learned from historical
data of existing centers. Each reusable model consists of a model architecture
and its source-trained parameter initialization and serves as a candidate for
target-center instantiation. Let the finite set of $N$ reusable source-trained
models be

\begin{equation}
\mathcal{Q}=\{q_n\}_{n=1}^{N},
\qquad
q_n=(a_n,\theta_n^0),
\label{eq:reusable-models}
\end{equation}

where $q_n$ denotes the $n$th reusable source-trained model,
$a_n\in\mathcal{A}$ denotes its architecture, and
$\theta_n^0\in\Theta(a_n)$ denotes its source-trained initialization. The
architectures represented in $\mathcal{Q}$ form the set
$\mathcal{A}_{\mathcal{Q}}=\{a_n\}_{n=1}^{N}\subseteq\mathcal{A}$.

For target center $c$, the platform specifies two per-model deployment
limits. They are a maximum estimated operation count for one batch-size-one
forward pass and a maximum number of trainable parameters. Let $B_c^{F}$ and $B_c^{P}$
denote these limits, respectively. For $a\in\mathcal{A}_{\mathcal{Q}}$, let
$F_c(a)$ denote the estimated operation count under the fixed counting rule,
and let $P_c(a)$ denote the exact number of trainable parameters. In this
formulation, parameter adaptation changes only parameter values while keeping
the architecture fixed. Therefore, $F_c(a)$ and $P_c(a)$ remain unchanged.

For center $c$, the deployment-feasible reusable model set is

\begin{equation}
\mathcal{Q}_{c}^{\mathrm{fea}}
=
\left\{
q_n\in\mathcal{Q}
\;\middle|\;
F_c(a_n)\leq B_c^{F},
\;
P_c(a_n)\leq B_c^{P}
\right\}.
\label{eq:feasible-candidate-set}
\end{equation}

Both deployment limits are modeled as hard constraints rather than penalty
terms.

\subsection{DT Model Instantiation Problem}

Let $\mathbb{P}_c$ denote the target-center data distribution over
$\mathcal{X}\times\mathcal{Y}$, and let
$\ell:\mathcal{Y}\times\mathcal{Y}\rightarrow\mathbb{R}_{\geq 0}$ denote the
task loss function. The target risk is defined as

\begin{equation}
R_c(a,\theta)
=
\mathbb{E}_{(x,y)\sim\mathbb{P}_c}
\left[
\ell(f_{a,\theta}(x),y)
\right].
\label{eq:target-risk}
\end{equation}

Let $\mathcal{U}_c$ denote a prescribed target-adaptation operator. For each
reusable model $q_n=(a_n,\theta_n^0)$,
$\mathcal{U}_c(q_n,\mathcal{D}_c^{\mathrm{sup}})\in\Theta(a_n)$ denotes
the target-specific parameters obtained by adapting $q_n$ using the target
support set. An ideal reusable-model decision for center $c$ is defined as

\begin{equation}
\begin{aligned}
n_c^\dagger
\in
\operatorname*{arg\,min}_{n\in\{1,\ldots,N\}}
&\quad
R_c\!\left(
a_n,
\mathcal{U}_c(q_n,\mathcal{D}_c^{\mathrm{sup}})
\right)
\\
\mathrm{s.t.}\quad
&F_c(a_n)\leq B_c^F,
\\
&P_c(a_n)\leq B_c^P.
\end{aligned}
\label{eq:instantiation-problem}
\end{equation}

Here, $n_c^\dagger$ denotes the reusable-model index selected under the
constrained target-risk criterion in Eq.~\eqref{eq:instantiation-problem}. The
corresponding DT model is
$\left(a_c^\dagger,\theta_c^\dagger\right)
=
\left(
a_{n_c^\dagger},
\mathcal{U}_c(q_{n_c^\dagger},\mathcal{D}_c^{\mathrm{sup}})
\right)$.
The target distribution $\mathbb{P}_c$ is unavailable during instantiation.
The target risk $R_c$ therefore cannot be evaluated directly. The decision
must instead be based on the support and validation data in
Eq.~\eqref{eq:target-data}.

\section{The Proposed MSA-DTI Algorithm}
\label{sec:method}

\subsection{Overview}

Figure~\ref{fig:method-overview} summarizes MSA-DTI for the instantiation
problem in Eq.~\eqref{eq:instantiation-problem}. The method contains an offline
preparation stage and a target-side instantiation stage. Offline preparation
constructs the reusable candidate bank $\mathcal{Q}$, designates the reference
candidate $q^{\mathrm{ref}}$, and calibrates the relative margin $\tau$ on
development data that are disjoint from final target evaluation. For a target
center, the target-side stage filters candidates by the two deployment limits,
adapts the feasible candidates using the same support set and update budget,
and compares their post-adaptation performance on the validation set. The best
alternative replaces the adapted reference only when it satisfies the
calibrated relative-improvement requirement. The output is the target-specific
DT model $\mathrm{DT}_c$.

\enlargethispage{8pt}
\vspace{-8pt}
\paperfig[trim={60bp 500bp 20bp 80bp},clip]{0.9\linewidth}
{fig2.pdf}
{Overall workflow of MSA-DTI.}
{fig:method-overview}

\subsection{Candidate Bank Construction}

Before constructing the final source-trained candidate bank, MSA-DTI performs
architecture preselection to identify the architectures to retain. This step uses a
dedicated pool of development target-center cases that is disjoint from the
data used for margin calibration and final target evaluation. A reference
architecture $a^{\mathrm{ref}}\in\mathcal{A}$ is fixed before preselection
and retained regardless of the preselection outcome. Each preselection case
contains three disjoint subsets. The support subset is used for adaptation, the
validation subset is used to select an alternative architecture, and the check
subset provides preselection evidence. For
each preselection case, the reference architecture must satisfy the assigned
deployment limits, and only feasible alternative architectures are included in
the comparison. During preselection, the reference and each feasible
alternative are adapted from fixed source-trained initializations on the
support subset before validation and check losses are evaluated.

For each alternative architecture, two evidence counts are accumulated across
the preselection cases. A \emph{check-oracle win} is assigned to the feasible
alternative architecture with the lowest check-set loss if its loss is lower
than that of the adapted reference. A
\emph{validation-selected positive check win} is assigned to the feasible
alternative architecture selected on the validation subset if its check-set
loss is lower than that of the adapted reference. The architecture is retained
if at least one of the two counts reaches a preset threshold.

Let $\mathcal{D}_{\mathrm{src}}$ denote the combined historical data from the
existing centers. For each retained architecture $a_q$, source pretraining on
$\mathcal{D}_{\mathrm{src}}$ produces an initialization $\theta_q^0$, forming
a candidate $q=(a_q,\theta_q^0)\in\mathcal{Q}$. For a retained architecture
with multiple source-trained initializations, each initialization defines a
distinct candidate. One candidate associated with the reference architecture
$a^{\mathrm{ref}}$ is designated as $q^{\mathrm{ref}}$.
Figure~\ref{fig:source-bank} summarizes the retained architectures and
resulting candidate bank. The candidate bank and reference candidate are fixed
before margin calibration and final target evaluation. This stage outputs
$\mathcal{Q}$ and $q^{\mathrm{ref}}$ for target-side instantiation.

\paperfig[trim={200bp 1200bp 60bp 140bp},clip]{0.92\linewidth}
{fig3.pdf}
{Candidate bank construction and reference candidate designation.}
{fig:source-bank}

\subsection{Target-Side Instantiation}

For target center $c$, MSA-DTI computes $\mathcal{Q}_c^{\mathrm{fea}}$ using
Eq.~\eqref{eq:feasible-candidate-set} and restricts target adaptation to the
candidates in this set. Because the subsequent reference-based selection
requires a deployment-feasible reference for comparison and fallback, MSA-DTI
proceeds only if $q^{\mathrm{ref}}\in\mathcal{Q}_c^{\mathrm{fea}}$.
Otherwise, it terminates without DT instantiation under the current deployment
limits.

For target-center instantiation, the support set
$\mathcal{D}_c^{\mathrm{sup}}$ is used for parameter adaptation, whereas the
validation set $\mathcal{D}_c^{\mathrm{val}}$ is used for post-adaptation
candidate comparison. For a dataset $\mathcal{D}$, let
\begin{equation}
\mathcal{L}(\mathcal{D};a,\theta)
=
\frac{1}{|\mathcal{D}|}
\sum_{(x,y)\in\mathcal{D}}
\ell\!\left(f_{a,\theta}(x),y\right)
\label{eq:empirical-loss}
\end{equation}
denote the empirical task loss.

For each feasible candidate $q=(a_q,\theta_q^0)$, let
$\theta_{c,q}^{(t)}$ denote its parameters after $t$ target-side adaptation
steps. With a common update budget of $T$ steps and a learning rate $\eta$
shared by all feasible candidates, the gradient update is

\begin{equation}
\begin{aligned}
\theta_{c,q}^{(0)}
&=
\theta_q^{0},
\\
\theta_{c,q}^{(t+1)}
&=
\theta_{c,q}^{(t)}
-
\eta
\nabla_{\theta}
\mathcal{L}
\left(
\mathcal{D}_{c}^{\mathrm{sup}};
a_q,
\theta_{c,q}^{(t)}
\right),
\\[-1mm]
&\hspace{38mm}
t=0,\ldots,T-1.
\end{aligned}
\label{eq:target-update}
\end{equation}

Equation~\eqref{eq:target-update} gives the unclipped form of the
target-adaptation operator $\mathcal{U}_c$ introduced in Section~3.3.
The experiments apply the gradient-norm limit in Table~\ref{tab:configuration}
before each update. The final parameters $\theta_{c,q}^{(T)}$ are the
operator's output. After $T$ adaptation steps, let $J_c(q)$ denote the
post-adaptation validation score of candidate $q$.

\begin{equation}
J_c(q)
=
\mathcal{L}
\left(
\mathcal{D}_{c}^{\mathrm{val}};
a_q,
\theta_{c,q}^{(T)}
\right).
\label{eq:validation-loss}
\end{equation}

Because validation-based candidate rankings are subject to estimation
uncertainty, MSA-DTI uses a relative margin so that a marginal post-adaptation
loss reduction does not by itself trigger replacement of the reference
candidate.

Let
$\mathcal{Q}_{c}^{\mathrm{alt}}
=
\mathcal{Q}_{c}^{\mathrm{fea}}
\setminus
\{q^{\mathrm{ref}}\}$
denote the feasible alternatives, and let $q_c^*$ denote the candidate
selected by MSA-DTI. If $\mathcal{Q}_{c}^{\mathrm{alt}}$ is empty, MSA-DTI
sets $q_c^*=q^{\mathrm{ref}}$ and skips alternative selection. Otherwise, let
$\widetilde q_c
\in
\operatorname*{arg\,min}_{q\in\mathcal{Q}_{c}^{\mathrm{alt}}}
J_c(q)$
denote an alternative with the lowest validation loss. Let $\tau\in[0,1)$
denote the relative margin. The selected candidate is then

\begin{equation}
q_c^{*}
=
\begin{cases}
\widetilde q_c,
&
J_c(\widetilde q_c)
\leq
(1-\tau)J_c(q^{\mathrm{ref}}),
\\
q^{\mathrm{ref}},
&
\text{otherwise}.
\end{cases}
\label{eq:reference-based-selection}
\end{equation}

A larger $\tau$ requires a larger validation-loss reduction for an
alternative to be selected over the reference. The margin affects only model
selection and does not change the adapted parameters.

The selected candidate $q_c^*$ provides the architecture
$a_{q_c^*}$, and its adapted parameters $\theta_{c,q_c^*}^{(T)}$ define the
target-specific DT model in Eq.~\eqref{eq:dt-instance}. This stage outputs
$\mathrm{DT}_c$. Figure~\ref{fig:adaptation-selection} illustrates
common-budget adaptation and the reference-based selection rule.

\paperfig[trim={140bp 690bp 100bp 120bp},clip]{0.88\linewidth}
{fig5.pdf}
{Common-budget adaptation and reference-based selection.}
{fig:adaptation-selection}

Algorithm~\ref{alg:msa-dti} summarizes the target-side procedure when a
calibrated margin is available. The reference-only fallback when calibration
fails is specified in Section~4.4.

\begin{algorithm}[H]
\caption{Target-Side MSA-DTI for One Center}
\label{alg:msa-dti}
\begin{algorithmic}[1]

\Require Candidate bank $\mathcal{Q}$, reference candidate
$q^{\mathrm{ref}}$, support set $\mathcal{D}_{c}^{\mathrm{sup}}$,
validation set $\mathcal{D}_{c}^{\mathrm{val}}$, deployment limits
$B_c^{F}$ and $B_c^{P}$, update budget $T$,
learning rate $\eta$, and relative margin $\tau$

\Ensure Target-specific DT model $\mathrm{DT}_c$, if the reference candidate is deployment-feasible

\State Compute $\mathcal{Q}_{c}^{\mathrm{fea}}$ using
Eq.~\eqref{eq:feasible-candidate-set}

\If{$q^{\mathrm{ref}}\notin\mathcal{Q}_{c}^{\mathrm{fea}}$}
    \State terminate without DT instantiation
\EndIf

\ForAll{$q\in\mathcal{Q}_{c}^{\mathrm{fea}}$}
    \State Adapt $q$ using Eq.~\eqref{eq:target-update}
    \State Compute $J_c(q)$ using Eq.~\eqref{eq:validation-loss}
\EndFor

\State
$\mathcal{Q}_{c}^{\mathrm{alt}}
\gets
\mathcal{Q}_{c}^{\mathrm{fea}}
\setminus
\{q^{\mathrm{ref}}\}$

\If{$\mathcal{Q}_{c}^{\mathrm{alt}}=\emptyset$}
    \State $q_c^{*}\gets q^{\mathrm{ref}}$
\Else
    \State Choose
    $\widetilde q_c
    \in
    \underset{q\in\mathcal{Q}_{c}^{\mathrm{alt}}}
    {\arg\min}\;J_c(q)$
    \State Select $q_c^{*}$ using
    Eq.~\eqref{eq:reference-based-selection}
\EndIf

\State
$\mathrm{DT}_{c}
\gets
\left(
a_{q_c^{*}},
\theta_{c,q_c^{*}}^{(T)}
\right)$

\State \Return $\mathrm{DT}_{c}$

\end{algorithmic}
\end{algorithm}

\subsection{Relative-Margin Calibration and Analytical Properties}

The margin $\tau$ is calibrated on a margin-calibration pool that is disjoint
from the architecture-preselection cases and final target evaluation. Each
calibration case contains disjoint support, validation, and check sets. The
support and validation sets retain the roles defined in Section~4.3, whereas the check set
provides independent evidence for evaluating the resulting selection during
calibration. For a given margin, the harmful-selection rate is the fraction of
all calibration cases in which an alternative is selected and its check-set
loss exceeds that of the adapted reference. The paired-initialization control
uses a second source-trained initialization of the reference architecture and
provides a same-architecture control for source-initialization effects.

Let $\mathcal{T}$ be a preset margin set. For $\tau\in\mathcal{T}$, let
$\widehat R_{\mathrm{harm}}(\tau)$ denote the observed harmful-selection rate.
Let $\widehat G_{\mathrm{ref}}(\tau)$ and
$\widehat G_{\mathrm{pair}}(\tau)$ denote the mean relative check-loss gains
with respect to the adapted reference and paired-initialization control,
respectively, where each gain is the relative reduction in check-set loss of
the selected model with respect to the corresponding control. Their lower
confidence bounds are $\underline G_{\mathrm{ref}}(\tau)$ and
$\underline G_{\mathrm{pair}}(\tau)$. Given a harmful-selection-rate
threshold $\rho$ and mean-gain thresholds $\gamma_{\mathrm{ref}}$ and
$\gamma_{\mathrm{pair}}$, the eligible margins are

\begin{equation}
\begin{aligned}
\mathcal{T}_{\mathrm{elig}}
=
\bigl\{
\tau\in\mathcal{T}
\;\bigm|\;&
\widehat R_{\mathrm{harm}}(\tau)\leq\rho,
\\
&
\widehat G_{\mathrm{ref}}(\tau)\geq\gamma_{\mathrm{ref}},
\quad
\underline G_{\mathrm{ref}}(\tau)>0,
\\
&
\widehat G_{\mathrm{pair}}(\tau)\geq\gamma_{\mathrm{pair}},
\quad
\underline G_{\mathrm{pair}}(\tau)>0
\bigr\}.
\end{aligned}
\label{eq:margin-calibration}
\end{equation}

When $\mathcal{T}_{\mathrm{elig}}$ is nonempty, the calibrated margin is the
least restrictive eligible margin $\tau^*=\min\mathcal{T}_{\mathrm{elig}}$
and is fixed before final target evaluation.

If $\mathcal{T}_{\mathrm{elig}}=\varnothing$, MSA-DTI retains the adapted
reference. Because the statistics in Eq.~\eqref{eq:margin-calibration} are
estimated from the calibration pool, the calibrated margin depends on the
calibration domain.

By construction, any DT model instantiated by MSA-DTI is selected from
$\mathcal{Q}_{c}^{\mathrm{fea}}$ and therefore satisfies both deployment
limits. Moreover, whenever an alternative is selected and
$J_c(q^{\mathrm{ref}})>0$, Eq.~\eqref{eq:reference-based-selection}
guarantees a relative validation-loss reduction of at least $\tau$ with
respect to the reference.

The target-stage computational complexity is
\begin{equation}
O\left(
|\mathcal{Q}|+
\sum_{q\in\mathcal{Q}_{c}^{\mathrm{fea}}}
\left(TC_q^{\mathrm{grad}}+C_q^{\mathrm{val}}\right)
\right).
\label{eq:target-complexity}
\end{equation}
Here, $C_q^{\mathrm{grad}}$ and $C_q^{\mathrm{val}}$ denote the costs of one
target update and validation evaluation for candidate $q$, respectively. The
$|\mathcal{Q}|$ term accounts for deployment filtering.

\section{Experimental Evaluation}
\label{sec:experiments}

We evaluate prediction accuracy, selection reliability, deployment
feasibility, target-side cost, and robustness across target conditions.

\subsection{Experimental Setup}
\label{subsec:experimental-setup}

\subsubsection{Benchmark and Data Splits}

We use a controlled synthetic benchmark for heterogeneous multi-center
workloads. The benchmark contains three center types A--C with distinct
workload-dynamics ranges. Monitoring becomes progressively sparser and noisier
from type A to type C. The prediction target is the next $H$ values of a
normalized latent workload in $[0,1]$; 12 simulated monitoring metrics
with value/mask pairs and a time-gap feature constitute the 25-channel input.
Synthetic sampling intervals are measured in simulated base steps rather than
an asserted physical clock interval; MSE is evaluated on the squared workload
fraction scale. Supplementary Section~S8 gives the type-specific generation
ranges and missingness rules. Deployment-limit tiers are platform-side settings
assigned using the type-dependent probabilities specified in Supplementary
Section~S8.

The main study uses 20 source centers and 20 held-out target centers.
Prediction horizons $H\in\{1,4\}$ and support sizes $K\in\{10,20\}$ give
four cases per target center and 80 held-out cases in total. Here, $K$
denotes the number of support windows used for target adaptation.

For each case, the raw timeline is partitioned chronologically into disjoint
support, validation, check, and test segments before window generation, and
windows are constructed separately within each segment. Depending on $H$ and
$K$, the validation, check, and test sets contain 1342--1351, 621--628, and
621--627 windows, respectively. Across the architecture-preselection, margin-calibration, and diagnostic pools,
check sets provide the corresponding development evidence, whereas held-out
target test sets are reserved for final evaluation after model selection. The source, architecture-preselection,
margin-calibration, diagnostic, and held-out target pools use disjoint center
identities.

\subsubsection{Candidate Bank, Reference, and Calibration}
\label{subsubsec:bank-calibration}

The initial architecture space contains 66 configurations spanning multilayer
perceptron (MLP), temporal convolutional network (TCN), and gated recurrent
unit (GRU) architectures. A57, a one-layer GRU with 32 hidden units, is fixed
as the reference before preselection.

An alternative configuration is retained if either of the two preselection
evidence counts defined in Section~4.2 is at least two. The preselection step retains six configurations corresponding to five executable
architectures. Together with a second source-trained initialization of A57 used
as the paired-initialization control, the resulting source-trained records form
a seven-candidate bank. The reference, paired A57 control, and stronger
alternative initializations do not share a single identical source-training
recipe; therefore, effects of the complete bank and selector must not be
interpreted as architecture-only effects.

Applying Eq.~\eqref{eq:margin-calibration} to the six preset margins on
20 disjoint calibration centers (80 cases) yields the smallest eligible
margin $\tau=0.10$, which is fixed before held-out evaluation. The full
five-criterion calibration grid is reported in Supplementary Section~S2.2.
Table~\ref{tab:configuration} summarizes the experimental configuration.

\begin{center}
\begin{minipage}{\linewidth}
\captionof{table}{Main experimental configuration.}
\label{tab:configuration}
\centering
\scriptsize
\renewcommand{\arraystretch}{1.05}
\setlength{\tabcolsep}{2.5pt}

\begin{tabularx}{\linewidth}{Z{0.18\linewidth} Y Z{0.43\linewidth}}
\toprule
Category & Item & Setting \\
\midrule

Data
& Centers, observations, and input length
& Source/held-out $=20/20$, 3000 per center, $L=96$ \\

& Target cases and split sizes
& 80 cases, $H\in\{1,4\}$, support $K\in\{10,20\}$,
validation/check/test windows
$=1342$--$1351$ / $621$--$628$ / $621$--$627$ \\

\midrule

Bank
& Initial / retained / executable / candidates
& $66/6/5/7$ \\

& Reference and architecture preselection
& A57, one-layer GRU-32, fixed before preselection. An alternative is retained if
either preselection evidence count in Section~4.2 is $\geq 2$ \\

& Source-model records
& Reference pooled-source initialization, five alternative source-trained
initializations, and one additional A57 paired-initialization control \\

\midrule

Training
& Strong-bank alternative source training
& Adam / MSE / 50 epochs, learning rate $10^{-3}$, batch size 64 \\

& MSA-DTI and fixed-reference baseline adaptation
& SGD / MSE / 50 updates per adapted model, learning rate 0.01,
gradient-norm limit 1.0, full support set \\

\midrule

Selection
& Grid and eligibility
& $\{0.05, 0.075, 0.10, 0.125, 0.15, 0.20\}$,
harmful-selection rate $\leq 5\%$, both mean relative gains $\geq 3\%$,
both 95\% lower confidence bounds $>0$ \\

& Calibrated margin
& $\tau=0.10$ \\

\midrule

Limits
& Operation-count limits
& Tight / medium / loose
$=1.5{\times}10^6/5{\times}10^6/2{\times}10^7$ \\

& Parameter-count limits
& Tight / medium / loose
$=3{\times}10^4/10^5/5{\times}10^5$ \\

\midrule

Statistics
& Confidence intervals
& Center-cluster bootstrap, 4000 repetitions \\

\bottomrule
\end{tabularx}
\end{minipage}
\end{center}

\subsubsection{Task, Baselines, and Metrics}

We use workload prediction as the evaluation task because it supports
proactive resource management in computing systems
\cite{ref17,ref18,ref19}. Mean squared error (MSE) is the task loss. We also
report mean absolute error (MAE), Worst-10\% error, conditional value at risk
at the 90th percentile (CVaR90), target-side time, and limit satisfaction.

For each case, Worst-10\% error is the mean of the largest 10\% of window-level
MSE values on the test set. CVaR90 is the mean case-level MSE over the 10\% of
target cases with the largest MSE.

For the main held-out cases, beneficial and harmful alternative selections
are reported against the independently evaluated pre-training and fine-tuning
(PT+FT) baseline using test MSE. This differs from the within-run adapted
reference used for the calibration check-set comparison. Relative gains are
computed for each paired case and then averaged. The 95\% confidence intervals
use a center-cluster bootstrap with 4000 repetitions, resampling target centers
while keeping their four cases together.

We compare MSA-DTI with PT+FT, MeDeT-based adaptation, random
initialization, three task-matched NAS-inspired controls (abbreviated as
few-shot NAS, zero-shot NAS, and zero-shot NAS with fine-tuning in figures
and tables), and H-Meta-NAS-based adaptation \cite{ref14,ref15,ref30}.
These controls use the common architecture space with method-specific
shortlisting and target evaluation; the abbreviations do not denote exact
reproductions of the published NAS algorithms. PT+FT is the main
controlled baseline because it uses the
same source-trained reference and nominal target-adaptation settings as
MSA-DTI. The remaining baselines provide method-level comparisons while retaining their
method-specific initialization or search procedures.

As the closest evaluated architecture-adaptation baseline, H-Meta-NAS-based
adaptation uses source-learned meta-initializations and validation-only target
search after deployment filtering. Its meta-initializations and search settings are fixed before held-out
evaluation, and the main comparison uses 50 SGD/MSE updates per adapted
model. All methods use the same target
cases, chronological splits, deployment limits, and architecture definitions.

\subsubsection{Runtime Protocol}

All methods are implemented in Python and PyTorch and run on an NVIDIA RTX
3060 Laptop GPU. Target-side time includes model creation, initialization
loading, target-side search when applicable, target adaptation, validation
evaluation, and final selection, but excludes offline preparation and test
evaluation. Each reported time is the mean over five runs.

\subsection{Overall Performance and Selection Reliability}
\label{subsec:overall}

Figure~\ref{fig:overall-performance} compares predictive performance on the
held-out target cases. All selected models satisfy both deployment limits.

\paperfig[trim={0bp 10bp 0bp 0bp},clip]{0.98\linewidth}
{fig_overall_performance_ours.pdf}
{Predictive performance on the held-out target cases. Panels (a)--(d) show MAE, MSE, Worst-10\% error, and CVaR90, respectively. Lower values are better.}
{fig:overall-performance}

MSA-DTI gives the lowest values for all four error measures. Compared with
PT+FT, the strongest baseline, MSA-DTI reduces MAE by 8.55\%, MSE by 14.60\%,
and Worst-10\% error by 13.76\% on average across paired cases. The 95\%
confidence interval for the paired MSE reduction is $[9.41\%,20.14\%]$.
Paired MSE comparisons with the remaining baselines also favor MSA-DTI, with
each 95\% confidence interval lying above zero.

Figure~\ref{fig:paired-instantiation} compares MSA-DTI with PT+FT at the case
level.

\paperfig{0.70\linewidth}
{fig6.pdf}
{Held-out paired comparison of MSA-DTI and PT+FT for (a) test MSE and (b) test Worst-10\% error. Alternative labels use test MSE against the independently evaluated PT+FT baseline, and points below the diagonal favor MSA-DTI.}
{fig:paired-instantiation}

MSA-DTI retains the reference in 33 cases and selects an alternative in 47
cases. Of the 47 alternative selections, 44 have lower test MSE than the
independently evaluated PT+FT baseline and three have higher test MSE.
These three represent 3.75\% of all held-out cases. The 95\% confidence
interval for this PT+FT-based harmful-alternative rate is
$[0.00\%,7.50\%]$, which spans 5\% and does not directly validate the
within-run calibration criterion.

Because few-shot NAS and zero-shot NAS with fine-tuning use Adam and Huber
loss in the main comparison, we additionally evaluate an optimizer-matched
diagnostic in which the included methods use SGD, MSE, and 50 target updates
per adapted model. In this
diagnostic, MSA-DTI achieves the lowest MSE, Worst-10\% error, and CVaR90 among
the included methods. Table~\ref{tab:optimizer-matched} reports these results.

\begin{center}
\begin{minipage}{\linewidth}
\captionof{table}{Optimizer-matched diagnostic comparison.}
\label{tab:optimizer-matched}
\centering
\scriptsize
\renewcommand{\arraystretch}{1.07}
\setlength{\tabcolsep}{2.8pt}
\begin{tabularx}{\linewidth}
{Y Z{0.17\linewidth} Z{0.20\linewidth} Z{0.18\linewidth} Z{0.18\linewidth}}
\toprule
Method
& MSE$\downarrow$
& Worst-10\%$\downarrow$
& CVaR90$\downarrow$
& Adapted models per case \\
\midrule
PT+FT & 0.004497 & 0.017144 & 0.011878 & 1.00 \\
Few-shot NAS & 0.006715 & 0.025246 & 0.019653 & 12.00 \\
Zero-shot NAS + fine-tuning & 0.007402 & 0.028099 & 0.020402 & 12.00 \\
\textbf{MSA-DTI} & \textbf{0.003809} & \textbf{0.015094} &
\textbf{0.008696} & 6.10 \\
\bottomrule
\end{tabularx}
\vspace{2pt}
\par\footnotesize
All methods use the same independent diagnostic cases, hard deployment
limits, SGD optimizer, MSE loss, and 50 target updates per adapted model. The
source initializations and candidate-selection rules remain method-specific.
\end{minipage}
\end{center}

\subsection{Robustness and Sensitivity Analysis}
\label{subsec:heterogeneity}

Figure~\ref{fig:heterogeneity-robustness} reports paired MSE reductions
relative to PT+FT. Positive values favor MSA-DTI.

\paperfig{0.88\linewidth}
{fig7.pdf}
{Paired MSE reduction relative to PT+FT across (a) horizon/support settings and (b) target-center types. Error bars show 95\% confidence intervals.}
{fig:heterogeneity-robustness}

The reduction ranges from 11.57\% to 18.03\% across the four horizon and
support-size settings and from 13.03\% to 16.80\% across the three center
types. All seven confidence intervals remain above zero.

On a separate robustness pool, the fixed candidate architectures retrained
under three source-training seeds yield paired MSE reductions ranging from
10.52\% to 13.21\%. All three 95\% confidence intervals remain above zero.

Applying the same full 66-configuration architecture-preselection procedure to
two new 20-center pools yields preselection-derived lists of zero and four
indexed configurations, respectively, compared with six in the main-study
list. The fixed A57 reference candidate is constructed separately and remains
available for downstream fallback regardless of these list sizes. The two independent
retained sets have Jaccard similarities of 0 and 0.429 to the main-study
retained set. A separate architecture-preselection stability diagnostic evaluates the two
nonempty retained sets on a fourth independent 20-center pool while keeping
the downstream procedure unchanged. The 95\% confidence intervals for the
paired check-MSE gains relative to the reference include zero for both
nonempty sets.

The retained-set variation indicates sample sensitivity in candidate
discovery. To examine how this variability propagates through the complete
instantiation procedure, we then construct the corresponding source-trained
banks, apply the same margin-calibration rule to each bank, and evaluate
the resulting pipelines on a new 20-center held-out pool. The banks derived
from the main-study preselection pool and the nonempty independent preselection pool
achieve mean paired MSE reductions of 10.03\%
(95\% CI $[4.60\%,16.02\%]$) and 8.78\%
(95\% CI $[3.83\%,14.24\%]$) relative to PT+FT, respectively. Each produces
six harmful alternative selections, corresponding to 7.50\% of the 80
held-out cases. The preselection outcome with no retained alternative therefore
reverts to the adapted reference, yielding zero paired difference. All three
pipelines satisfy both deployment limits in every case.

\subsection{Deployment Feasibility and Cost}
\label{subsec:deployment}

Feasibility in this subsection is assessed using the operation-count and
parameter-count limits defined in Section~3.2. Figure~\ref{fig:budget-architecture}
shows how the three deployment-limit settings affect candidate feasibility and
final model selection.

\paperfig{0.98\linewidth}
{fig8.pdf}
{Effects of the three deployment-limit settings on
(a) the number of feasible candidates and
(b) the selected model configurations.}
{fig:budget-architecture}

Under the tight setting, four of the seven candidates in the main candidate
bank are deployment-feasible, whereas all seven are feasible under the medium and loose
settings.

For the main candidate bank under the tight setting, a single-limit analysis
shows that the reduction is caused by the parameter-count limit. All seven
candidates satisfy the operation-count limit, while the parameter-count limit
excludes three MLP candidates.

To test whether the two limits independently restrict candidate
feasibility, we form an eight-candidate stress bank by adding one configuration
from the predefined architecture space and use diagnostic thresholds derived
from the resulting model-complexity profiles. For both $H=1$ and $H=4$, each
single-limit feasible set retains seven of eight candidates, whereas the joint
feasible set retains six.

Under the joint stress limits, all models selected from the expanded bank on
the same 80 held-out cases satisfy both limits. Separately, for the main
candidate bank, mean paired MSE reductions relative to PT+FT remain positive
under the tight, medium, and loose settings. The corresponding reductions
are 11.35\%, 15.30\%, and 18.22\%, respectively.

Table~\ref{tab:target-cost} compares target-side cost and selected-model
complexity for PT+FT, the target-adaptation architecture-search baselines, and
MSA-DTI.

\begin{center}
\begin{minipage}{\linewidth}
\captionof{table}{Target-side cost and selected-model complexity.}
\label{tab:target-cost}
\centering
\scriptsize
\renewcommand{\arraystretch}{1.07}
\setlength{\tabcolsep}{2.8pt}
\begin{tabularx}{\linewidth}
{Y Z{0.20\linewidth} Z{0.18\linewidth} Z{0.16\linewidth} Z{0.19\linewidth}}
\toprule
Method
& Time per case (s)
& Adapted models per case
& Parameter count
& Estimated operation count \\
\midrule
PT+FT
& $0.204\pm0.004$
& 1.00
& 5{,}746
& 1{,}050{,}784 \\
Few-shot NAS
& $13.471\pm0.173$
& 12.00
& 30{,}995
& 1{,}468{,}894 \\
H-Meta-NAS
& $20.847\pm5.032$
& 12.00
& 26{,}627
& 1{,}700{,}279 \\
Zero-shot NAS + fine-tuning
& $17.601\pm0.150$
& 12.00
& 62{,}076
& 1{,}571{,}621 \\
\textbf{MSA-DTI}
& $5.676\pm0.059$
& 6.25
& 14{,}892
& 751{,}433 \\
\bottomrule
\end{tabularx}
\vspace{2pt}
\par\footnotesize
Time is mean $\pm$ standard deviation over five runs. Parameter
counts and estimated operation counts are averaged over the final models
selected in the held-out cases.
\end{minipage}
\end{center}


Relative to PT+FT, MSA-DTI achieves lower held-out MSE in the main
evaluation but takes 5.676 rather than 0.204 seconds per target case
(approximately 27.8 times, or 5.472 seconds more). It is \(2.37\times\)
faster than few-shot NAS and \(3.10\times\) faster than zero-shot NAS
with fine-tuning. These comparisons concern one-time target-side model
instantiation, not concurrent platform scheduling. The
reported target-side times are measured for one instantiation case on one GPU.
Concurrent-case scaling and platform-level scheduling are outside the present
evaluation.

Hardware profiling evaluates every executable architecture on an Intel Core
i7-12700H CPU and an NVIDIA RTX 3060 Laptop GPU. Median batch-size-one
inference latency ranges from 0.093 to 1.980 ms on the CPU and from 0.133 to
4.170 ms on the GPU. Serialized model size ranges from 10.3 to 315.9 KiB.
Operation count is positively associated with median latency, whereas
parameter count preserves the serialized-size ranking in all measured
settings. The measured increment in peak PyTorch CUDA allocation is not
monotonic in parameter count. These measurements support the two counts as
distinct model-level complexity indicators. Their hardware-level
implications remain host-dependent.

\subsection{Ablation Study}
\label{subsec:component}

Table~\ref{tab:component} reports component variants on an independent
diagnostic pool. The reference-only and two-reference-initialization variants
do not reproduce the cross-architecture gain. Without the relative margin,
diagnostic MSE is nearly unchanged, while the harmful-selection rate increases
from 15.00\% to 17.50\%. The individual rate intervals overlap, but this is
not a paired test of the difference. The increase is reported descriptively. Removing deployment
filtering slightly lowers MSE but reduces limit satisfaction from 100\% to
96.25\%.

Overall, the main held-out result evaluates the complete MSA-DTI pipeline,
while the diagnostic controls distinguish the roles of individual design
choices. The cross-architecture gain remains sensitive to source training. The
relative margin mainly affects replacement conservativeness, whereas
deployment filtering enforces limit satisfaction.

\begin{center}
\begin{minipage}{\linewidth}
\captionof{table}{Component variants on the independent diagnostic pool.}
\label{tab:component}
\centering
\scriptsize
\renewcommand{\arraystretch}{1.08}
\setlength{\tabcolsep}{2.2pt}
\begin{tabularx}{\linewidth}
{Y Z{0.12\linewidth} Z{0.14\linewidth} Z{0.17\linewidth} Z{0.16\linewidth} Z{0.16\linewidth}}
\toprule
Variant
& MSE$\downarrow$
& Worst-10\%$\downarrow$
& Paired MSE reduction (\%)$\uparrow$
& Harmful-selection rate (\%)$\downarrow$
& Limit satisfaction (\%)$\uparrow$ \\
\midrule
\textbf{MSA-DTI} & 0.004063 & 0.015764 & 15.79 & 15.00 & 100.00 \\
Earlier source training & 0.004876 & 0.019044 & 1.75 & 2.50 & 100.00 \\
Reference only & 0.005051 & 0.019512 & 0.00 & 0.00 & 100.00 \\
Two reference initializations & 0.004917 & 0.019136 & 1.29 & 1.25 & 100.00 \\
Without relative margin & 0.004055 & 0.015690 & 16.16 & 17.50 & 100.00 \\
Without deployment filtering & 0.004037 & 0.015635 & 16.13 & 13.75 & 96.25 \\
\bottomrule
\end{tabularx}
\vspace{2pt}
\par\footnotesize
Paired MSE reduction is computed relative to the adapted reference on the same
diagnostic pool. Harmful-selection rate is computed over all diagnostic
cases. The earlier source-training row changes the full source-training
procedure for the alternatives, so it measures candidate-bank sensitivity
rather than a single-component ablation.
\end{minipage}
\end{center}

\subsection{Real-Workload Evaluations}
\label{subsec:external}

The primary real-workload experiment uses the GWA-T-12 Bitbrains
fastStorage trace\footnote{Dataset: \url{https://atlarge-research.com/gwa-t-12/}.
The trace was provided by Bitbrains IT Services Inc.}
of virtual-machine resource measurements \cite{ref37}.
All source models are trained anew within this data domain, using 20 source
VMs; separate 20-VM development and calibration pools precede a 40-VM
held-out pool (160 $H$/$K$ cases). This is an in-domain real-data
instantiation study, not zero-shot cross-dataset transfer. It predicts
future recorded CPU usage from the next $H\in\{1,4\}$ native records
using $K\in\{10,20\}$ support windows; native record gaps vary, with
nominal five-minute sampling. Under the full five-criterion calibration
rule, $\tau=0.05$ is selected before held-out testing. The mean of
case-relative test-MSE reductions against PT+FT is 32.70\%
(95\% VM-cluster CI [7.67\%,51.25\%]); against a candidate fixed by
independent development VMs under the same feasibility limits, the
secondary paired reduction is 38.39\% (95\% CI [25.56\%,50.80\%]).
These ratios differ from the 8.09\% reduction computed using the two
aggregate mean MSE values. Eleven of 160 held-out cases undergo a harmful
alternative selection (6.875\%), exceeding the 5\% empirical calibration
criterion, which does not provide a held-out risk guarantee.
The multi-candidate method also incurs additional target-side adaptation
cost. Detailed dataset, calibration, cost, and subgroup records are provided
in Supplementary Section~S7.5. A separate synthetic strong-fixed control
with unresolved historical development-pool provenance is reported as
supplementary evidence in Section~S7.6 rather than a confirmatory main result.

To examine an additional workload setting, we evaluate Alibaba Cluster
Trace v2018\footnote{Alibaba Cluster Trace v2018 is available at
\url{https://github.com/alibaba/clusterdata/tree/master/cluster-trace-v2018}.}
under a corrected causal-input protocol. Source models are retrained on
past-only features; 20 source, 20 calibration, and 40 held-out machines
yield 80 calibration and 160 held-out cases. Each model predicts $H\in\{1,4\}$ future CPU-fraction values
from 96 historical steps on a 10-second event-time grid, using
$K\in\{10,20\}$ support windows for adaptation. Observed CPU values,
rather than interpolated future labels, provide the test targets. No preset margin
passes the full five-criterion calibration, so MSA-DTI retains the
adapted PT+FT reference throughout the held-out evaluation (0.00\%
paired gain). Applying the synthetic-study margin $\tau=0.10$ only as
a diagnostic yields $-0.2592\%$ mean case-relative test-MSE gain
(95\% machine-cluster CI $[-1.4711\%,0.6205\%]$), without
an established positive gain on this trace. This retrospective correction involves
previously observed machines and assumes event timestamps represent
data-availability times. Supplementary Section~S7.4 provides the
complete protocol, calibration table, and held-out analyses.

\section{Conclusion}
\label{sec:conclusion}

This paper presented MSA-DTI for DT model instantiation with few-shot target
adaptation under per-model operation-count and parameter-count limits. MSA-DTI filters
infeasible candidates from the source-trained bank, adapts the remaining
candidates under a common update budget, and applies a calibrated
reference-model replacement rule based on post-adaptation validation
performance. In the main held-out evaluation, MSA-DTI reduces MSE by 14.60\%
on average across paired cases relative to PT+FT while satisfying both
deployment limits in all main held-out cases. The confidence interval for the
PT+FT-based held-out harmful-alternative rate spans 5\%. Candidate
discovery also shows sensitivity to the development pool used for architecture
preselection. In a separate end-to-end analysis, the two nonempty
preselection-derived banks yield positive mean paired MSE reductions relative
to PT+FT with confidence intervals above zero, while the no-alternative
preselection outcome reverts to the adapted reference. In a separately source-trained and calibrated Bitbrains VM cohort,
MSA-DTI achieves a 32.70\% mean case-relative test-MSE reduction over
PT+FT; the same-bank development-fixed comparison provides secondary
support. The method incurs additional target-side adaptation cost, and
11 of 160 held-out Bitbrains cases have harmful alternative selections,
above the empirical 5\% calibration criterion. In the causally corrected
Alibaba evaluation, no preset margin qualifies, so the calibrated policy
falls back to the adapted reference. These findings support positive
performance on the evaluated synthetic and Bitbrains workloads without
implying universal gains across real-data domains. Future work includes
more stable candidate discovery, selection under domain shift, and
measurement-derived deployment limits across additional platforms.

\section*{CRediT Authorship Contribution Statement}

\textit{The final author contribution statement will be added after the
author list is confirmed.}

\section*{Funding}

This work is supported by the Major Science and Technology Projects in Qinghai
Province (2024-GX-A3).

\section*{Declaration of Competing Interest}

The authors declare that they have no known competing financial
interests or personal relationships that could have appeared to
influence the work reported in this paper.

\section*{Data Availability}

The frozen synthetic-study code, configurations, processed results,
plot-ready data, and reproduction scripts are available in GitHub Release
v1.2.7 at
\url{https://github.com/w924871681/paper-code-DT-Twin/releases/tag/v1.2.7}.
The expanded R5 and Bitbrains assets include code, configurations, model
checkpoints, saved predictions and targets, development-set fixed-candidate
mappings, calibration and held-out records, statistical verification scripts,
and reproduction instructions. Corrected Alibaba assets include preprocessing
and evaluation code, checkpoints, six-threshold calibration records,
held-out results, and reproduction instructions. These assets are available at
\url{https://github.com/w924871681/paper-code-DT-Twin/releases/tag/evidence-full-20261013};
the earlier light-audit archive remains available at
\url{https://github.com/w924871681/paper-code-DT-Twin/releases/tag/evidence-r5-bitbrains-20261013}.
The archives exclude third-party raw traces and their full derived input
windows. R5's historical R2 development-pool RNG/raw-array provenance remains
\texttt{UNRESOLVED\_HISTORICAL\_RNG\_PROVENANCE}. Corrected Alibaba is a causal
reconstruction on a previously observed evaluation population, not a new blind
evaluation. Independent statistical recalculation of saved results does not
establish independent end-to-end training reproduction from raw data. The original Alibaba Cluster Trace
v2018 and Bitbrains GWA-T-12 traces are not redistributed; their
official access links are provided in the corresponding dataset footnotes.

\section*{Declaration of Generative AI and AI-Assisted Technologies
in the Manuscript Preparation Process}

During the preparation of this work, the authors used OpenAI ChatGPT to
assist with language refinement, manuscript organization, and readability
improvement. After using this tool, the authors reviewed and edited the
content as needed and take full responsibility for the content of the
published article.

\begin{thebibliography}{00}

\bibitem{ref1} L. U. Khan, Z. Han, W. Saad, E. Hossain, M. Guizani, and C. S. Hong, ``Digital twin of wireless systems: Overview, taxonomy, challenges, and opportunities,'' \emph{IEEE Commun. Surveys Tuts.}, vol. 24, no. 4, pp. 2230--2254, 4th Quart. 2022, doi: 10.1109/COMST.2022.3198273.

\bibitem{ref2} S. Mihai et al., ``Digital twins: A survey on enabling technologies, challenges, trends and future prospects,'' \emph{IEEE Commun. Surveys Tuts.}, vol. 24, no. 4, pp. 2255--2291, 4th Quart. 2022, doi: 10.1109/COMST.2022.3208773.

\bibitem{ref3} H. Xu, J. Wu, Q. Pan, X. Guan, and M. Guizani, ``A survey on digital twin for Industrial Internet of Things: Applications, technologies and tools,'' \emph{IEEE Commun. Surveys Tuts.}, vol. 25, no. 4, pp. 2569--2598, 4th Quart. 2023, doi: 10.1109/COMST.2023.3297395.

\bibitem{ref4} F. Tao, H. Zhang, A. Liu, and A. Y. C. Nee, ``Digital twin in industry: State-of-the-art,'' \emph{IEEE Trans. Ind. Informat.}, vol. 15, no. 4, pp. 2405--2415, 2019, doi: 10.1109/TII.2018.2873186.

\bibitem{ref8} J. Li et al., ``Digital twin-enabled service provisioning in edge computing via continual learning,'' \emph{IEEE Trans. Mobile Comput.}, vol. 23, no. 6, pp. 7335--7350, 2024, doi: 10.1109/TMC.2023.3332668.

\bibitem{ref11} H. Sartaj, S. Ali, and J. M. Gj{\o}by, ``MeDeT: Medical device digital twins creation with few-shot meta-learning,'' \emph{ACM Trans. Softw. Eng. Methodol.}, vol. 34, no. 6, Art. no. 158, pp. 1--36, 2025, doi: 10.1145/3708534.

\bibitem{ref9} A. Vettoruzzo, M.-R. Bouguelia, J. Vanschoren, T. R{\"o}gnvaldsson, and K. C. Santosh, ``Advances and challenges in meta-learning: A technical review,'' \emph{IEEE Trans. Pattern Anal. Mach. Intell.}, vol. 46, no. 7, pp. 4763--4779, 2024, doi: 10.1109/TPAMI.2024.3357847.

\bibitem{ref10} H. Gharoun, F. Momenifar, F. Chen, and A. H. Gandomi, ``Meta-learning approaches for few-shot learning: A survey of recent advances,'' \emph{ACM Comput. Surv.}, vol. 56, no. 12, Art. no. 294, pp. 1--41, 2024, doi: 10.1145/3659943.

\bibitem{ref12} M. Tan et al., ``MnasNet: Platform-aware neural architecture search for mobile,'' in \emph{Proc. IEEE/CVF Conf. Comput. Vis. Pattern Recognit. (CVPR)}, 2019, pp. 2820--2828, doi: 10.1109/CVPR.2019.00293.

\bibitem{ref13} B. Wu et al., ``FBNet: Hardware-aware efficient ConvNet design via differentiable neural architecture search,'' in \emph{Proc. IEEE/CVF Conf. Comput. Vis. Pattern Recognit. (CVPR)}, 2019, pp. 10726--10734, doi: 10.1109/CVPR.2019.01099.

\bibitem{ref26} Y. Zhao, L. Wang, and T. Guo, ``Multi-objective neural architecture search by learning search space partitions,'' \emph{J. Mach. Learn. Res.}, vol. 25, no. 177, pp. 1--41, 2024.

\bibitem{ref30}
Y. Zhao, X. Gao, I. Shumailov, N. Fusi, and R. Mullins,
``Rapid model architecture adaption for meta-learning,''
in \emph{Adv. Neural Inf. Process. Syst.},
vol. 35, 2022, pp. 18721--18732,
doi: 10.52202/068431-1360.

\bibitem{ref31}
D. Lian et al.,
``Towards fast adaptation of neural architectures with meta learning,''
in \emph{Proc. Int. Conf. Learn. Represent. (ICLR)}, 2020.

\bibitem{ref32}
P. Eustratiadis, L. Dudziak, D. Li, and T. Hospedales,
``Neural fine-tuning search for few-shot learning,''
in \emph{Proc. Int. Conf. Learn. Represent. (ICLR)}, 2024.

\bibitem{ref5} B. Qin et al., ``Machine and deep learning for digital twin networks: A survey,'' \emph{IEEE Internet Things J.}, vol. 11, no. 21, pp. 34694--34716, Nov. 2024, doi: 10.1109/JIOT.2024.3416733.

\bibitem{ref6} A. Hakiri, A. Gokhale, S. Ben Yahia, and N. Mellouli, ``A comprehensive survey on digital twin for future networks and emerging Internet of Things industry,'' \emph{Comput. Netw.}, vol. 244, Art. no. 110350, 2024, doi: 10.1016/j.comnet.2024.110350.

\bibitem{ref7} Y. Wu, K. Zhang, and Y. Zhang, ``Digital twin networks: A survey,'' \emph{IEEE Internet Things J.}, vol. 8, no. 18, pp. 13789--13804, Sep. 2021, doi: 10.1109/JIOT.2021.3079510.

\bibitem{ref27}
M. Emami Khansari and S. Sharifian,
``A deep reinforcement learning approach towards distributed Function as a
Service (FaaS) based edge application orchestration in cloud-edge
continuum,''
\emph{J. Netw. Comput. Appl.}, vol. 233,
Art. no. 104042, 2025,
doi: 10.1016/j.jnca.2024.104042.

\bibitem{ref28} D. Hortelano et al.,
``A comprehensive survey on reinforcement-learning-based computation
offloading techniques in Edge Computing Systems,''
\emph{J. Netw. Comput. Appl.}, vol. 216,
Art. no. 103669, 2023, doi: 10.1016/j.jnca.2023.103669.

\bibitem{ref24} J. Son, S. Lee, and G. Kim, ``When meta-learning meets online and continual learning: A survey,'' \emph{IEEE Trans. Pattern Anal. Mach. Intell.}, vol. 47, no. 1, pp. 413--432, 2025, doi: 10.1109/TPAMI.2024.3463709.

\bibitem{ref20} L. M. Zintgraf, K. Shiarlis, V. Kurin, K. Hofmann, and S. Whiteson, ``Fast context adaptation via meta-learning,'' in \emph{Proc. 36th Int. Conf. Mach. Learn. (ICML)}, vol. 97, 2019, pp. 7693--7702.

\bibitem{ref21} Y. Zhao, T. Zhang, J. Li, and Y. Tian, ``Dual adaptive representation alignment for cross-domain few-shot learning,'' \emph{IEEE Trans. Pattern Anal. Mach. Intell.}, vol. 45, no. 10, pp. 11720--11732, 2023, doi: 10.1109/TPAMI.2023.3272697.

\bibitem{ref22} J. Wu, X. Liu, X. Yin, T. Zhang, and Y. Zhang, ``Task-adaptive prompted transformer for cross-domain few-shot learning,'' in \emph{Proc. AAAI Conf. Artif. Intell.}, vol. 38, no. 6, pp. 6012--6020, 2024, doi: 10.1609/aaai.v38i6.28416.

\bibitem{ref23} W. Wang, L. Duan, Y. Wang, J. Fan, and Z. Zhang, ``MMT: Cross domain few-shot learning via meta-memory transfer,'' \emph{IEEE Trans. Pattern Anal. Mach. Intell.}, vol. 45, no. 12, pp. 15018--15035, Dec. 2023, doi: 10.1109/TPAMI.2023.3306352.

\bibitem{ref25} Z. Zuo, Y. Huang, Z. Li, Y. Jiang, and C. Liu, ``Mixed contrastive transfer learning for few-shot workload prediction in the cloud,'' \emph{Computing}, vol. 107, no. 1, Art. no. 5, 2025, doi: 10.1007/s00607-024-01366-y.

\bibitem{ref33}
J. Bornschein, F. Visin, and S. Osindero,
``Small data, big decisions: Model selection in the small-data regime,''
in \emph{Proc. 37th Int. Conf. Mach. Learn. (ICML)},
vol. 119, 2020, pp. 1035--1044.

\bibitem{ref34}
A. Lacoste, M. Marchand, F. Laviolette, and H. Larochelle,
``Agnostic Bayesian learning of ensembles,''
in \emph{Proc. 31st Int. Conf. Mach. Learn. (ICML)},
vol. 32, 2014, pp. 611--619.

\bibitem{ref35}
K. You, Y. Liu, J. Wang, and M. Long,
``LogME: Practical assessment of pre-trained models for transfer learning,''
in \emph{Proc. 38th Int. Conf. Mach. Learn. (ICML)},
vol. 139, 2021, pp. 12133--12143.

\bibitem{ref36}
K. You, Y. Liu, Z. Zhang, J. Wang, M. I. Jordan, and M. Long,
``Ranking and tuning pre-trained models: A new paradigm for exploiting model hubs,''
\emph{J. Mach. Learn. Res.}, vol. 23, no. 209, pp. 1--47, 2022.

\bibitem{ref14} Y. Zhao, L. Wang, Y. Tian, R. Fonseca, and T. Guo, ``Few-shot neural architecture search,'' in \emph{Proc. 38th Int. Conf. Mach. Learn. (ICML)}, vol. 139, 2021, pp. 12707--12718.

\bibitem{ref15} G. Li, D. Hoang, K. Bhardwaj, M. Lin, Z. Wang, and R. Marculescu, ``Zero-shot neural architecture search: Challenges, solutions, and opportunities,'' \emph{IEEE Trans. Pattern Anal. Mach. Intell.}, vol. 46, no. 12, pp. 7618--7635, Dec. 2024, doi: 10.1109/TPAMI.2024.3395423.

\bibitem{ref16} M.-T. Wu, H.-I. Lin, and C.-W. Tsai, ``A training-free neural architecture search algorithm based on search economics,'' \emph{IEEE Trans. Evol. Comput.}, vol. 28, no. 2, pp. 445--459, Apr. 2024, doi: 10.1109/TEVC.2023.3264533.

\bibitem{ref17} Y.-M. Kim, S. Song, B.-M. Koo, J. Son, Y. Lee, and J.-G. Baek, ``Enhancing long-term cloud workload forecasting framework: Anomaly handling and ensemble learning in multivariate time series,'' \emph{IEEE Trans. Cloud Comput.}, vol. 12, no. 2, pp. 789--799, Apr.--Jun. 2024, doi: 10.1109/TCC.2024.3400859.

\bibitem{ref18} F. Zhao, W. Lin, S. Lin, S. Tang, and K. Li, ``MSCNet: Multi-scale network with convolutions for long-term cloud workload prediction,'' \emph{IEEE Trans. Serv. Comput.}, vol. 18, no. 2, pp. 969--982, Mar.--Apr. 2025, doi: 10.1109/TSC.2025.3536313.

\bibitem{ref19} B. Feng and Z. Ding, ``Application-oriented cloud workload prediction: A survey and new perspectives,'' \emph{Tsinghua Sci. Technol.}, vol. 30, no. 1, pp. 34--54, Feb. 2025, doi: 10.26599/TST.2024.9010024.

\bibitem{ref37}
S. Shen, V. van Beek, and A. Iosup,
``Statistical characterization of business-critical workloads hosted in cloud datacenters,''
in \emph{Proc. 15th IEEE/ACM Int. Symp. Cluster, Cloud and Grid Comput. (CCGrid)},
2015, pp. 465--474, doi: 10.1109/CCGrid.2015.60.

\end{thebibliography}

\end{document}
